\scanpage{13}
\textcircled{III}\quad Polygônes $n$-réguliers et repr. du \uline{groupe} \inserted{diédral}

\deleted{\uline{Revenons au cas des polygônes réguliers}.}
\deleted{\unsure{S'il|Il|Soit} est $\Pi$ épinglé, \inserted{où $\Pi=(S,\omega=\{u,u^{-1}\})$,} et si \unsure{$\Gamma$ est le|$\Gamma$ est $k$}}
% slanted
\margin{\deleted{Question} Parler de l'épinglage (ou \uline{repère}) du pt de vue combinatoire (revient à la donnée, soit d'un couple $(s_0,d_0)$ d'un sommet et d'une \deleted{\unsure{arête}} \unsure{de|et}, soit au couple de deux sommets consécutifs $(s_0,s_1)$, soit au couple $(s_0,u)$ d'un sommet et d'une orientation. Donner la propriété de \uline{rigidification} (1)}
\[
\deleted{\Gamma=\operatorname{Aut}(\Pi)\qquad
\Gamma^{\unsure{\bullet|0}}=\text{ens des autom conservant l'orient.}}
\]
\deleted{on obtient une représentation}
\[
\deleted{\Gamma\xrightarrow{\varphi}\operatorname{Aut\,aff}(E)}
\]
\deleted{et $S$ \inserted{ensemble des sommets} est donc une orbite \unsure{ss|sous} $\Gamma$.}

\deleted{\unsure{du}} \inserted{Nous partirons du} \uline{pt de vue intrinsèque} \deleted{(non épinglé)}.

Le polygône $P=(S,\omega)$ détermine un
sous-groupe
\[
\Gamma\subset\operatorname{Aut\,aff}(E)
\]
isomorphe à un groupe diédral $D_n$
\[
\Gamma\simeq D_n\quad \inserted{n\in\mathbb N\ \text{convenable}}\quad
\text{(isomorphisme \uline{non} donné)}\tag*{(1)}
\]
et $S$ est une orbite de $\Gamma$ op. sur $E$
\[
S\in\deleted{\unsure{E/\Gamma|\Gamma}}\Gamma\backslash E\tag*{(2)}
\]
donc le ss-gpe $\Gamma$ étant donné, \deleted{$\unsure{E|S}$}
$S$ est déterminé quand on connaît
\uline{un} de ses sommets, soit
\[
\unsure{\cdot}\ s_0\in E\tag*{(4)}
\]
\deleted{\unsure{si que|à qu|À quelle}} Le polygône \unsure{$P$|$P_s$} est-il connu
par la connaissance de $(\Gamma,\deleted{s_0,}S)$,
i.e. la structure combinatoire \inserted{$\omega$} de $S$
résulte-t-elle de la connaissance du
ss-gpe $\Gamma\subset\deleted{\operatorname{Aut}}\unsure{\mathfrak S|\circlearrowright}_S$ de ses autom.?
% archive sheet number: 12

